% !TeX program = xelatex
% Main manuscript: Elsevier CAS single-column, Chinese working draft.
% Original CAS sample and the previous manuscript are preserved in archive/.
\providecommand{\PaperRoot}{./}
\makeatletter
\edef\input@path{{\PaperRoot}{\PaperRoot template/}}
\makeatother
\PassOptionsToPackage{no-math}{fontspec}
\documentclass[a4paper,fleqn]{cas-sc}
\usepackage[UTF8,scheme=plain,fontset=windows]{ctex}
\usepackage[numbers,sort&compress]{natbib}
\usepackage{bm,tabularx,longtable}
\usepackage{algorithm,algpseudocode}
\usepackage{tikz}
\usetikzlibrary{arrows.meta,positioning,fit}
\usepackage{placeins,needspace}
\hypersetup{unicode=true,hidelinks}
\pdfstringdefDisableCommands{\def\\{ }}
\graphicspath{{\PaperRoot}}
\numberwithin{equation}{section}
\numberwithin{figure}{section}
\numberwithin{table}{section}
\renewcommand{\figurename}{图}
\renewcommand{\tablename}{表}
\renewcommand{\abstractname}{摘要}
\floatname{algorithm}{算法}
\algrenewcommand{\algorithmicrequire}{\textbf{输入：}}
\renewcommand{\refname}{参考文献}
\newcommand{\R}{\mathbb R}
\newcommand{\C}{\mathbb C}
\newcommand{\E}{\mathbb E}
\newcommand{\Rea}{\operatorname{Re}}
\newcommand{\tr}{\operatorname{tr}}
\newcommand{\diag}{\operatorname{diag}}
\setlength{\emergencystretch}{2em}
\widowpenalty=10000
\clubpenalty=10000
% No author information has been supplied. Keep the draft footer factual.
\ExplSyntaxOn
\cs_set:Npn \__first_footerline:
  { \small\rmfamily\itshape Signal~Processing\quad 中文初稿 }
% CAS 2.4 uses a zero-width keyword overlay; stretchable glue preserves its
% position while avoiding the template's artificial overfull-box diagnostic.
\exp_args:Nc \patchcmd { environment~Abstract~code }
  { \box \g_stm_key_box } { \box \g_stm_key_box \hss }
  { } { \PackageError{paper}{CAS~abstract~compatibility~patch~failed}{} }
\cs_new_eq:NN \paper_printorcid: \printorcid
\RenewDocumentCommand \printorcid { }
  { \seq_if_empty:NF \g_stm_orcid_seq { \paper_printorcid: } }
\ExplSyntaxOff
\begin{document}
\let\WriteBookmarks\relax
\shorttitle{未知宽带强弱重叠源的高效直接定位}
\shortauthors{}
\title[mode=title]{未知宽带强弱重叠源的高效直接定位：\\多候选联合精修与波形消元更新}
% Add the real author, affiliation and correspondence here before submission.

% BEGIN sections/abstract
\begin{abstract}[摘要]
两步定位先估计站间到达时间差，再由时差求解位置；当未知宽带源占用同一频带、接收功率相差较大时，弱源时延提取和跨站关联容易受到强源干扰。直接位置估计（DPD）通过联合拟合多站观测处理这一问题，其计算效率和恢复能力取决于候选搜索与连续更新：粗谱排名可能遗漏精修后更优的解，固定波形系数的位置更新又可能产生大量小步。本文在自由空间相对幅度和理想基带时延模型下，研究同步单天线站的双源定位，形成多候选联合精修AP（MCJ-AP）和波形消元块NOMP（VP-BNOMP）。MCJ-AP保留少量条件候选，对每个分支同时优化两源位置，再按完整拟合残差择优；VP-BNOMP在每次局部位置试探时重估逐频信号系数，以消元曲率处理位置与波形的耦合。在20个源站布局、六档功率差的2400条观测中，五候选MCJ-AP将双源20\,m成功率由单候选的95\%提高至99\%，平均耗时0.718\,s。相同五候选、相同停止精度的配对表明，联合精修保持定位表现并减少31.36\%的完整耗时。共同240条观测中，VP-BNOMP与固定系数参照的成功数相同，平均耗时减少79.92\%。候选消融、曲率分析和受控更新对照分别说明了可靠性与计算收益的来源。
\end{abstract}
\begin{keywords}[关键词]
直接位置估计 \sep 强弱重叠源 \sep 交替投影 \sep 块NOMP \sep 可分离非线性最小二乘
\end{keywords}

% END sections/abstract

\maketitle

% BEGIN sections/introduction
\section{引言}\label{sec:introduction}
辐射源定位常先估计到达时间差等中间量，再由这些量求解位置。广义互相关与双曲线定位为这一流程提供了成熟工具\cite{KnappCarter1976,ChanHo1994}。当多个源使用重叠频带且波形未知时，各站接收的是不同传播响应下的混合信号，中间量的提取与跨站关联随之变难。直接位置估计（direct position determination，DPD）将传播关系直接带入观测模型，以位置作为主要待估参数，为多源联合处理提供了另一条途径\cite{ChenHudsonYao2002,WeissAmar2005}。

本文研究四个同步单天线站对两个未知宽带源的直接定位。两源占用相同频带，接收功率可以明显不同；定位器已知站点坐标及相对传播规律，不知道发射波形与发射功率。目标是在有限计算时间内准确恢复两个位置。单源空间谱中的强源及其旁瓣可能占据多个高分位置，直接取两个峰容易漏掉弱源；同时搜索两个二维位置则需要检查大量位置组合。逐源发现后再连续精修可以减少搜索量，但其效果取决于两个具体选择：保留哪些起点，以及怎样更新源位置。

已有直接定位方法从不同角度处理混合观测。宽带最大似然与交替投影（alternating projection，AP）固定部分位置、更新其余位置，降低联合搜索的维数\cite{ChenHudsonYao2002}。多重信号分类（MUSIC）型DPD根据候选响应与噪声子空间的正交性定位\cite{WeissAmar2005}；最小方差无失真响应（MVDR）利用协方差逆矩阵压低干扰\cite{Capon1969}，DPD-HR将这一思想用于直接定位\cite{TirerWeiss2016}。本文让这些方法使用相同观测、处理频带和传播信息，比较它们在强弱源混合时的实际表现。

稀疏参数估计将粗搜索与离网精修结合起来。Newtonized orthogonal matching pursuit（NOMP）在贪心发现后进行Newton更新和循环精修\cite{Mamandipoor2016}。Yang等的多频率二维块NOMP以距离衰减和传播相位构造响应，将各频率的未知源系数组成一个块，并在位置更新后重估系数\cite{Yang2022}。其宽带模型与本文接近，是局部更新的直接参照。QNOMP则针对已知导频形成的信道观测，先以最小二乘消去路径增益，再联合优化各路径的时延和角度\cite{Zeng2025}。本文沿用这些多频拟合与连续优化基础，集中改进两个环节：在条件搜索后保留多个可精修分支，以及在块NOMP的位置试探中同步消去当前源的逐频系数。

候选保留与后续拟合择优已有一般性研究。前瞻正交最小二乘（look-ahead orthogonal least squares，LAOLS）保留多个离散原子候选，向前补全支持集，再按最终残差决定当前选择\cite{Chatterjee2012}。本文面对的差别在于位置可连续变化：粗条件谱固定第一源位置，连续精修却允许两源共同移动，两种排序因而可能不同。MCJ-AP保留多个条件候选，分别与第一源组成分支，联合优化后按完整残差择优。新增计算集中在少量局部分支，不需要为每个分支重复全域扫描。

未知波形的处理同样影响迭代效率。变量投影法先求最优线性系数，再优化非线性参数\cite{GolubPereyra1973}。这里线性系数是逐频波形，非线性参数是源位置。试探位置改变时重新估计波形，可以避免固定系数对位置步长的约束。本文据此构造波形消元块NOMP，并推导固定系数曲率与消元曲率的关系，解释局部更新为何可能需要不同的计算量。

本文的主要工作如下。
\begin{enumerate}
  \item 围绕同步单天线站的未知宽带强弱重叠源，明确观测、传播和接收功率定义，以宽带最小二乘消元为基础，将条件搜索、残差发现和位置精修置于统一的拟合框架中。
  \item 提出多候选联合精修AP，通过精修后的完整残差选择分支，降低粗搜索过早舍弃有效位置的风险；在相同候选和停止精度下，联合位置更新进一步减少分支处理成本。
  \item 构造波形消元块NOMP，保持原有发现与循环结构，以消元曲率替换固定系数曲率。结合局部误差传播、数值保护隔离和同输入对照，说明位置--波形耦合导致小步更新的机制及其计算改进。
\end{enumerate}

下文依次给出观测模型、求解算法与计算量分析，再在20个源站布局、六档功率差下检验定位表现。全部2400条观测用于核心方法及消融，预先指定的240条共同观测用于八类方法比较。

% END sections/introduction


% BEGIN sections/model
\section{观测模型与定位目标}\label{sec:model}
\subsection{同步单天线站的宽带观测}
在静态直达传播条件下，源位置决定了信号到各站的相对幅度与时延。设$L$个单天线站的已知位置为$\bm r_l\in\R^2$，$l=1,\ldots,L$；$K$个源的未知位置为$\bm p_k\in\Omega\subset\R^2$，$k=1,\ldots,K$，其中$\Omega$为二维搜索区域。各站时间、频率同步，接收链路相对幅相响应已标定；发射天线各向一致，接收天线具有相同响应并满足极化匹配。本文先研究已知源数的定位，实验取$L=4$、$K=2$。距离、时延与相对幅度分别为
\begin{equation}
 d_{lk}=\|\bm p_k-\bm r_l\|_2,\qquad
 \tau_{lk}=d_{lk}/c_0,\qquad \beta_{lk}=d_{lk}^{-1},
 \label{eq:distance}
\end{equation}
其中$c_0$为传播速度，$d_{lk}$和$\tau_{lk}$的单位分别为m和s。

幅度模型来自自由空间路径损耗。对于远场、极化匹配的理想各向同性收发天线，接收功率$P_{\mathrm r}$与发射功率$P_{\mathrm t}$满足\cite{ITURP5252024}
\begin{equation}
 \frac{P_{\mathrm r}}{P_{\mathrm t}}
 =\left(\frac{\lambda}{4\pi d}\right)^2,\qquad
 a_{\mathrm{amp}}(d)=\sqrt{\frac{P_{\mathrm r}}{P_{\mathrm t}}}
 =\frac{\lambda}{4\pi d} .\label{eq:freespace}
\end{equation}
这里$\lambda$为工作波长，$d$为距离，$a_{\mathrm{amp}}$为幅度比。式\eqref{eq:freespace}说明，功率按$1/d^2$衰减，幅度按$1/d$衰减。工作波长和天线响应固定时，将公共因子$\lambda/(4\pi)$并入未知源波形，就得到$\beta_{lk}=1/d_{lk}$。定位器据此利用站间相对幅度，而无需知道发射功率。本文采用自由空间幅度规律与理想基带纯时延响应：给定位置后，跨站相位关系由基带频率和传播时延确定。接收链路标定用于消除站端固定响应；基带纯时延关系是另行规定的传播假设，不由链路标定自动推出。仿真与定位器均使用这一模型，不引入射频载波传播相位或额外随机链路增益。

令$s_k(t)$为已包含公共幅度因子和未知发射功率的第$k$个源基带波形。第$l$站以采样率$f_s$采集$N$点，接收信号为
\begin{equation}
 x_l[n]=\sum_{k=1}^{K}\beta_{lk}s_k(n/f_s-\tau_{lk})+w_l[n],
 \quad n=0,\ldots,N-1 .\label{eq:time}
\end{equation}
其中$x_l[n]$为复数IQ样本，$n$为采样索引。源波形互相独立且未知，在仿真中采用独立、零均值的随机符号产生。噪声$w_l[n]\sim\mathcal{CN}(0,\sigma_t^2)$为圆对称复高斯白噪声，$\sigma_t^2=\E|w_l[n]|^2$，在站间及采样点间独立，并与源信号独立。给定候选位置后，其到各站的相对幅度与时延均由式\eqref{eq:distance}确定。

为利用宽带信号的时延信息，对每站记录作离散傅里叶变换（DFT）：
\begin{equation}
 X_l(f_m)=\frac1N\sum_{n=0}^{N-1}x_l[n]
 \exp(-\mathrm j2\pi mn/N),\qquad f_m=mf_s/N .\label{eq:dft}
\end{equation}
$m$采用移位后的有符号频率索引，$\mathrm j=\sqrt{-1}$。在DFT频率轴上保留公共处理频带内的频点，组成集合$\mathcal F$，其大小记为$F=|\mathcal F|$。以$\bm x_f=[X_1(f),\ldots,X_L(f)]^T\in\C^L$收集同一频点的各站观测，则宽带频域模型写为\cite{ChenHudsonYao2002}
\begin{align}
 \bm x_f&=\bm A_f(\bm P)\bm s_f+\bm n_f,\quad f\in\mathcal F,\label{eq:frequency}\\
 [\bm a_f(\bm p)]_l&=\frac{1}{d_l(\bm p)}
       \exp[-\mathrm j2\pi f d_l(\bm p)/c_0],\label{eq:response}\\
 \bm A_f(\bm P)&=[\bm a_f(\bm p_1),\ldots,\bm a_f(\bm p_K)],\quad
 \bm P=[\bm p_1,\ldots,\bm p_K].\nonumber
\end{align}
其中$\bm s_f=[S_1(f),\ldots,S_K(f)]^T\in\C^K$为未知源频域系数，$\bm n_f\in\C^L$为噪声，$\bm a_f(\bm p)\in\C^L$为候选位置的跨站响应，$\bm A_f(\bm P)\in\C^{L\times K}$为各源响应组成的矩阵，$\bm P\in\R^{2\times K}$收集源坐标。式\eqref{eq:response}的相位项来自时移的傅里叶变换性质，$f$表示基带频率。全频带需要逐频计算响应，不能以单一频率的导向向量代替。

上述频域模型采用共同观测窗的时延近似，忽略各站有限记录在窗边缘的差异。完整DFT的白噪声满足$\E[\bm n_f\bm n_f^H]=(\sigma_t^2/N)\bm I_L$。去除带外频点后，保留频点的噪声方差不变，DFT归一化仍使用原始点数$N$。所有方法采用这一处理，不对各站混合IQ作额外RMS归一化。

计算中也可使用单位范数响应$\bm q_f(\bm p)=\bm a_f(\bm p)/\|\bm a_f(\bm p)\|_2$。这一操作对同一位置的整列响应除以同一个数，站间相对幅度不变。逐频源系数自由估计时，列缩放可由相应系数补偿，因此不改变联合拟合的最小残差。离散扫描包含站点坐标时，为避免$d_l=0$的数值奇点，实现将不足1\,m的距离按1\,m计算；本文的导数推导适用于这一保护区外，仿真源位置均远离站点。

下文粗体小写和大写分别表示向量与矩阵，$\C$和$\R$分别表示复数域和实数域；$(\cdot)^T$、$(\cdot)^H$及$(\cdot)^\dagger$分别表示转置、共轭转置和Moore--Penrose伪逆，$\|\cdot\|_2$表示欧氏范数，$\bm I_d$表示$d$阶单位矩阵。未标阶数的单位矩阵由相乘矩阵的维度确定。

\subsection{带内信噪比与强弱源条件}
记第$k$个源单独在第$l$站产生的时间平均接收功率为$P_{lk}$，名义信号带宽为$B$。对白噪声，带宽为$B$的频带占采样频带的比例为$B/f_s$，故其带内期望功率为$\sigma_t^2B/f_s$。以各站接收功率的平均值定义带内信噪比，以两个源接收总功率之比定义功率差：
\begin{equation}
 \mathrm{SNR}_{k,\mathrm{in}}
 =10\log_{10}\frac{L^{-1}\sum_{l=1}^L P_{lk}}
 {\sigma_t^2B/f_s},\qquad
 \Delta=10\log_{10}\frac{\sum_l P_{l1}}{\sum_l P_{l2}} .
 \label{eq:snr}
\end{equation}
式中$\mathrm{SNR}_{k,\mathrm{in}}$和$\Delta$均以dB表示。信噪比按名义信号带宽$B$定义，与含保护带的实际处理带宽区分。实验固定源2的波形、功率和噪声，仅缩放源1；这样改变功率差时，弱源自身的SNR保持不变。$\Delta>0$时称源1为强源、源2为弱源；$\Delta=0$时保留这两个身份编号。各站的局部功率差由传播距离决定，不必等于四站平均功率差。

\subsection{未知波形消元与联合位置估计}
定位时将逐频未知波形视为确定性待估系数。根据式\eqref{eq:frequency}的独立、等方差复高斯噪声模型，负对数似然中与位置和波形有关的部分为$\sigma_f^{-2}\sum_f\|\bm x_f-\bm A_f\bm s_f\|_2^2$，其中$\sigma_f^2=\sigma_t^2/N$。因此，最大似然估计等价于联合最小二乘估计\cite{ChenHudsonYao2002}。取观测总能量$E_x=\sum_{f\in\mathcal F}\|\bm x_f\|_2^2>0$作固定归一化，定义
\begin{align}
 \mathcal L(\bm P,\{\bm s_f\})
   &=E_x^{-1}\sum_f\|\bm x_f-\bm A_f(\bm P)\bm s_f\|_2^2,\label{eq:fullLS}\\
 \widehat{\bm s}_f(\bm P)&=\bm A_f(\bm P)^\dagger\bm x_f,\label{eq:waveLS}\\
 J(\bm P)&=E_x^{-1}\sum_f\|\bm\Pi_{\bm A_f}^{\perp}\bm x_f\|_2^2,
 \quad \bm\Pi_{\bm A_f}^{\perp}=\bm I-\bm A_f\bm A_f^\dagger .\label{eq:profile}
\end{align}
这里所有求和均遍历$\mathcal F$。给定一组位置，式\eqref{eq:waveLS}直接求出最优波形；$\bm A_f\bm A_f^\dagger$提取这组位置能够解释的观测，$\bm\Pi_{\bm A_f}^{\perp}\in\C^{L\times L}$保留剩余部分。将最优波形代回式\eqref{eq:fullLS}，就得到只含位置变量的目标$J$，即波形消元\cite{GolubPereyra1973}。$J$为无量纲残差，值越小，位置组合对观测的拟合越好。

所求位置为$\widehat{\bm P}=\arg\min_{\bm P\in\Omega^K}J(\bm P)$。当$K=2$时，非线性搜索只需处理两个二维位置，共四个实数。直接优化$J$时，每次试探新位置都共同重估两源波形。后文的AP联合精修采用这一方式；VP-BNOMP则交替更新位置与波形，在局部试探中只消去当前源的系数，逐源循环结束后再共同重估全部系数。

单源最小二乘谱则为
\begin{equation}
 Q_0(\bm p)=E_x^{-1}\sum_f
 \frac{|\bm a_f(\bm p)^H\bm x_f|^2}{\|\bm a_f(\bm p)\|_2^2} .\label{eq:single}
\end{equation}
式中$|\cdot|$表示复数模。单源情况下有$J(\bm p)=1-Q_0(\bm p)$，故最大化$Q_0$与最小化$J$等价。双源情况下，从同一张单源谱取两个峰没有同时拟合两源，不能代替对位置组合$\bm P$的优化。后续算法以单源谱发现初值，再用多源拟合处理相互影响。

% END sections/model


% BEGIN sections/methods
\section{多候选联合AP与波形消元块NOMP}\label{sec:methods}
本章分别改进候选选择和局部更新。MCJ-AP保留多个条件起点，以联合精修后的残差择优；VP-BNOMP在试探位置时重新估计当前源的波形，减少固定系数带来的重复小步。两种方法都求解联合拟合问题：前者直接优化完整消元目标$J$，后者使用固定其他源贡献的局部块目标，并在循环后共同重估波形。
\subsection{条件投影搜索}
寻找第二源时，应先区分哪些观测已由第一源解释。给定已发现位置，以$\bm B_f\in\C^{L\times K_{\mathrm B}}$收集其响应，$K_{\mathrm B}$为已固定的源数；搜索第二源时$K_{\mathrm B}=1$。对观测与候选响应作相同正交投影，定义
\begin{equation}
 \bm r_f=\bm\Pi_{\bm B_f}^{\perp}\bm x_f,\qquad
 \bm u_f(\bm p)=\bm\Pi_{\bm B_f}^{\perp}\bm a_f(\bm p).
\end{equation}
其中$\bm\Pi_{\bm B_f}^{\perp}=\bm I_L-\bm B_f\bm B_f^\dagger$，$\bm r_f\in\C^L$为投影后的观测，$\bm u_f(\bm p)\in\C^L$为候选响应经过同一投影后的向量。对应的条件位置谱为
\begin{equation}
 Q_{\mathrm{AP}}(\bm p\mid\bm B)
 =E_x^{-1}\sum_f\frac{|\bm u_f(\bm p)^H\bm r_f|^2}
 {\|\bm u_f(\bm p)\|_2^2}.\label{eq:conditional}
\end{equation}
条件谱衡量新增位置带来的拟合改善。将$\bm a_f$分解为已有响应空间内的部分与正交部分$\bm u_f$。当$\bm u_f\ne0$时，新方向与已有响应正交，投影矩阵满足
\begin{equation}
 \bm\Pi_{[\bm B_f,\bm a_f]}^{\perp}
 =\bm\Pi_{\bm B_f}^{\perp}
 -\frac{\bm u_f\bm u_f^H}{\|\bm u_f\|^2} .
\end{equation}
对上式作观测的二次型并求和，便有
\begin{equation}
 J_{\mathrm B}-J([\bm B,\bm a(\bm p)])
 =\frac1{E_x}\sum_f\frac{|\bm u_f(\bm p)^H\bm r_f|^2}
 {\|\bm u_f(\bm p)\|^2}
 =Q_{\mathrm{AP}}(\bm p\mid\bm B),\label{eq:conditionalGain}
\end{equation}
其中$J_{\mathrm B}=E_x^{-1}\sum_f\|\bm r_f\|^2$为仅拟合已有源的残差，$J([\bm B,\bm a(\bm p)])$表示加入候选后、共同重估所有系数所得的残差。若$\bm u_f=0$，该候选没有提供新的响应方向，此频点贡献记为零。

式\eqref{eq:conditionalGain}将条件搜索与联合拟合连接起来：固定已有位置时，条件得分最高的候选使联合残差最小\cite{ChenHudsonYao2002}。这一等价关系要求观测和候选响应使用相同投影。

\subsection{多候选联合精修AP（MCJ-AP）}\label{sec:mcj}
MCJ-AP以AP条件投影生成第二源候选，随后转入多分支联合位置优化。AP在此承担条件初始化，联合精修完成分支内求解；两者的具体分工见算法\ref{alg:ap}。
粗网格排序回答的是哪个网格点当前得分更高，连续精修则允许两源一起移动；二者得到的最佳分支可能不同。MCJ-AP据此保留多个可能起点。先在公共网格$\mathcal G\subset\Omega$上最大化$Q_0$，作单源连续精修，得到$\bm p_1^{(0)}$。固定该位置后，计算第二源条件谱，提取$3\times3$邻域内的局部极大值，再按得分从高到低保留间隔不少于$\Delta_{\mathrm c}$的点，得到$M$个候选$\bm c_1,\ldots,\bm c_M\in\R^2$。本文取$\Delta_{\mathrm c}=30$\,m。每个候选分别与第一源组成初值，进入相同的联合精修：
\begin{equation}
 \widehat{\bm P}^{(m)}=\operatorname{LocalMin}_{\bm P\in\Omega^2}
 J(\bm P;[\bm p_1^{(0)},\bm c_m]),\qquad
 m^*=\arg\min_m J(\widehat{\bm P}^{(m)}).\label{eq:branch}
\end{equation}
其中$\operatorname{LocalMin}$表示在搜索边界内从给定初值局部最小化$J$，$\widehat{\bm P}^{(m)}$为第$m$个分支的结果。算法输出最终残差最小的分支$\widehat{\bm P}^{(m^*)}$。因此，候选质量由精修后的完整拟合决定，粗谱仅负责提供起点。

\begin{algorithm}[H]
\caption{多候选联合精修AP}\label{alg:ap}
\begin{algorithmic}[1]
\Require 观测$\{\bm x_f\}$，站点坐标，公共网格$\mathcal G$，候选数$M$
\State 最大化$Q_0$并单源精修，得到$\bm p_1^{(0)}$
\State 固定$\bm p_1^{(0)}$，以式\eqref{eq:conditional}扫描$\mathcal G$
\State 提取邻域极大值，按条件得分和间距$\Delta_{\mathrm c}$保留$M$个候选
\For{$m=1,\ldots,M$}
  \State 从$[\bm p_1^{(0)},\bm c_m]$联合优化$J$；每次试探按式\eqref{eq:waveLS}重估波形
  \State 保存该分支的目标、位置与停止状态
\EndFor
\State 按式\eqref{eq:branch}输出最小残差分支
\end{algorithmic}
\end{algorithm}

联合精修采用有界拟牛顿法L-BFGS-B同时优化两个位置\cite{Byrd1995}。将坐标按源排列为$\bm\theta=\operatorname{vec}(\bm P)\in\R^{2K}$，其第$i$个分量为$\theta_i$；联合拟合残差为$\bm e_f=\bm x_f-\bm A_f\widehat{\bm s}_f\in\C^L$。响应矩阵局部满列秩时，最小二乘解满足$\bm A_f^H\bm e_f=0$。这一正交条件使求导中含$\partial\widehat{\bm s}_f/\partial\theta_i$的项消失，因而位置梯度可直接写为\cite{GolubPereyra1973}
\begin{equation}
 \frac{\partial J}{\partial\theta_i}
 =-\frac{2}{E_x}\Rea\sum_f
 \bm e_f^H\frac{\partial\bm A_f}{\partial\theta_i}
 \widehat{\bm s}_f,\label{eq:gradient}
\end{equation}
式中$\Rea$表示取实部，$\partial\bm A_f/\partial\theta_i\in\C^{L\times K}$为响应矩阵对该坐标的导数。记$\bm v_l=(\bm p-\bm r_l)/d_l\in\R^2$为从第$l$站指向候选位置的单位方向，$d_l=\|\bm p-\bm r_l\|_2$，响应的空间梯度为
\begin{equation}
 \nabla_{\bm p}a_l(f,\bm p)
 =a_l(f,\bm p)\left(-\frac1{d_l}-\mathrm j\frac{2\pi f}{c_0}\right)\bm v_l.
 \label{eq:responseDerivative}
\end{equation}
这里$a_l(f,\bm p)=[\bm a_f(\bm p)]_l$。括号内两项分别来自距离幅度和传播时延的变化。联合优化每次重估波形，但计算式\eqref{eq:gradient}时不需要显式求波形估计的导数。

增加候选数时，保持已有分支的初值、求解器与停止条件不变。较大集合包含较小集合的全部精修结果，故选择后的最小残差不会增大；对应的位置误差由实验评价。$M=1$为单候选版本，$M>1$以额外局部计算换取更多有效起点。

\subsection{波形消元块NOMP（VP-BNOMP）}\label{sec:vpnomp}
NOMP在当前拟合残差中发现新成分，再循环修正已有参数\cite{Mamandipoor2016}。二维多频率块NOMP将同一位置的逐频系数视为一个参数块\cite{Yang2022}。本文保留这一发现和循环结构，改进局部位置更新：每试探一个位置，就先求该源的最优信号系数，再判断位置是否更好。

给定当前位置和共同估计的系数$\widehat s_{k,f}$，更新源$j$时先减去其他源的当前拟合贡献：
\begin{equation}
 \bm y_{j,f}=\bm x_f-\sum_{k\ne j}\bm a_f(\bm p_k)\widehat s_{k,f} .
 \label{eq:targetResidual}
\end{equation}
其中$\bm y_{j,f}\in\C^L$为保留源$j$贡献的残差观测，在一次局部位置优化内固定。对任意试探位置$\bm p$，其最优复数系数$\widehat q_f(\bm p)\in\C$与局部目标$\phi_j(\bm p)$分别为
\begin{align}
 \widehat q_f(\bm p)&=\frac{\bm a_f(\bm p)^H\bm y_{j,f}}
 {\|\bm a_f(\bm p)\|^2},\label{eq:localcoef}\\
 \phi_j(\bm p)&=\min_{\{q_f\}}E_x^{-1}\sum_f
 \|\bm y_{j,f}-\bm a_f(\bm p)q_f\|^2\nonumber\\
 &=E_x^{-1}\sum_f\left(\|\bm y_{j,f}\|^2-
 \frac{|\bm a_f(\bm p)^H\bm y_{j,f}|^2}{\|\bm a_f(\bm p)\|^2}\right).
 \label{eq:localprofile}
\end{align}
这里$q_f$为标量信号系数，与单位范数响应向量$\bm q_f$区分。式\eqref{eq:localprofile}表示每个试探位置都配上其最优波形。Newton更新使用该消元目标的梯度与Hessian（二阶导数矩阵）；回溯时也重新求系数，仅接受局部残差不增的试探。更新后同步替换该源的响应、系数和显式残差。逐源循环结束后，再共同重估全部已发现源的系数。这样保持块NOMP的发现与循环结构，局部更新的区别集中在是否消去波形。

发现新源时，记当前共同拟合残差为$\bm r_f=\bm x_f-\sum_k\bm a_f(\bm p_k)\widehat s_{k,f}$；尚未发现源时$\bm r_f=\bm x_f$。扫描的残差谱为
\begin{equation}
 Q_{\mathrm R}(\bm p)=E_x^{-1}\sum_f
 \frac{|\bm a_f(\bm p)^H\bm r_f|^2}{\|\bm a_f(\bm p)\|^2},
 \label{eq:residualScore}
\end{equation}
该谱评价候选对当前残差的解释能力。它用于发现初始位置，后续再以联合拟合修正所有已发现源。

\begin{algorithm}[H]
\caption{波形消元块NOMP}\label{alg:nomp}
\begin{algorithmic}[1]
\Require 观测，公共网格，源数$K$，循环与Newton停止参数
\State 初始化源集合为空，$\bm r_f\gets\bm x_f$
\For{$k=1,\ldots,K$}
  \State 最大化残差谱\eqref{eq:residualScore}，发现新位置
  \State 以式\eqref{eq:localprofile}作单源Newton精修，加入源集合并更新残差
  \Repeat
    \For{已发现的每个源$j$}
      \State 构造$\bm y_{j,f}$，用$\phi_j$作有界Newton更新
      \State 每个试探位置重估$q_f$，仅接受局部残差不增的更新，同步替换该源贡献
    \EndFor
  \Until{循环停止条件满足或达到安全上限}
  \State 按式\eqref{eq:waveLS}共同重估全部已发现源的系数，更新共同拟合残差
\EndFor
\State 输出已发现位置
\end{algorithmic}
\end{algorithm}

两种发现方式的区别在于如何处理已有源。共同LS拟合后，残差$\bm r_f$与已有响应的列空间正交，因此AP条件谱和NOMP残差谱的分子相同；分母则分别为$\|\bm\Pi_{\bm B_f}^{\perp}\bm a_f\|^2$和$\|\bm a_f\|^2$。AP还考虑候选与已有响应的相似程度，NOMP直接扫描当前残差，两者因而可能选中不同的初始区域。本文保留这两种发现机制。

\subsection{消元更新的曲率与下降性质}\label{sec:curvature}
下面说明为什么固定信号系数可能使位置更新缓慢。对于一次源位置更新，将二维坐标记为$\bm z\in\R^2$，将$F$个复信号系数的实部和虚部排列为$\bm u\in\R^{2F}$。对应的归一化局部损失记为$\mathcal L_{\mathrm{loc}}(\bm z,\bm u)$，给定位置的最小二乘系数记为$\widehat{\bm u}(\bm z)$，消元后损失为$\varphi(\bm z)=\mathcal L_{\mathrm{loc}}(\bm z,\widehat{\bm u}(\bm z))$，即式\eqref{eq:localprofile}。

假设响应二阶连续可微，且系数最小二乘解局部唯一。定义$H_{zz}\in\R^{2\times2}$、$H_{uu}\in\R^{2F\times2F}$和$H_{zu}\in\R^{2\times2F}$分别为$\mathcal L_{\mathrm{loc}}$的位置、系数及混合二阶导数，$H_{uz}=H_{zu}^T$，所有导数在$(\bm z,\widehat{\bm u}(\bm z))$处计算。最优系数满足$\nabla_u\mathcal L_{\mathrm{loc}}=0$，对位置求导可得
\begin{equation}
 H_{uz}+H_{uu}\frac{\partial\widehat{\bm u}}{\partial\bm z}=0,
 \qquad\frac{\partial\widehat{\bm u}}{\partial\bm z}
 =-H_{uu}^{-1}H_{uz}.\label{eq:coefficientDerivative}
\end{equation}
将这一关系代入$\varphi$的一、二阶导数，得到变量投影的梯度和曲率关系\cite{GolubPereyra1973}：
\begin{align}
 \nabla\varphi&=\nabla_z\mathcal L_{\mathrm{loc}},\label{eq:envelope}\\
 H_{\mathrm{VP}}&=H_{zz}-H_{zu}H_{uu}^{-1}H_{uz},\label{eq:schur}\\
 H_{zz}-H_{\mathrm{VP}}&\succeq0 .\label{eq:psd}
\end{align}
其中$H_{\mathrm{VP}}=\nabla^2\varphi\in\R^{2\times2}$，下标VP表示变量投影；$\succeq0$表示半正定。式\eqref{eq:psd}来自系数最小二乘子问题的$H_{uu}$正定，因此$H_{zu}H_{uu}^{-1}H_{uz}$为半正定矩阵。

式\eqref{eq:schur}解释了两种更新的步长差别。位置移动引起的部分响应变化，可以由重新估计波形补偿。固定系数曲率$H_{zz}$没有扣除这部分变化，因而在耦合较强的方向上可能过大；相同梯度除以较大的正曲率，就得到较小步长。消元曲率$H_{\mathrm{VP}}$扣除了位置--波形耦合项，使更新直接针对重估波形后的残差。

在搜索区域内部的局部极小点$\bm z^*$附近，令位置误差$\bm\epsilon=\bm z-\bm z^*$。当$H_{zz}$和$H_{\mathrm{VP}}$均正定，且采用完整Newton步时，用固定系数曲率更新消元梯度的一阶误差传播为
\begin{equation}
 \bm\epsilon^{+}\approx
 \left(\bm I_2-H_{zz}^{-1}H_{\mathrm{VP}}\right)\bm\epsilon .\label{eq:localrate}
\end{equation}
式中$\bm\epsilon^+$为更新后的位置误差，Hessian取局部极小点处的值。由式\eqref{eq:psd}，迭代矩阵的特征值落在$[0,1)$；当某方向的消元曲率远小于固定系数曲率时，该方向的误差收缩因子接近1，迭代便可能出现大量小步。采用精确消元Newton曲率可去掉这一局部线性误差项。实际算法还受回溯、步长限制和Hessian近似的影响，因此用目标与导数调用次数进一步评价计算效率。

波形消元与联合精修处理两种不同的相互影响：前者允许波形随位置调整，后者允许两个位置一起调整。对联合目标$J$，双源的$4\times4$ Hessian为
\begin{equation}
 H_{\mathrm{joint}}=
 \begin{bmatrix}H_{11}&H_{12}\\H_{21}&H_{22}\end{bmatrix} .
\end{equation}
其中$H_{ij}=\partial^2J/(\partial\bm p_i\partial\bm p_j^T)\in\R^{2\times2}$。逐源精修每次只改变一个位置，联合精修同时调整两个位置，从而利用跨源项$H_{12}$和$H_{21}$所反映的相互影响。MCJ-AP采用联合L-BFGS-B更新；VP-BNOMP改进每个源内部的位置--系数耦合，不额外叠加末端联合优化。

MCJ-AP的联合精修仅接受完整目标不增的更新，即$J(\bm P^+)\le J(\bm P)$。VP-BNOMP的每次块更新降低或保持显式拟合残差，循环后的共同最小二乘重估也不增加该残差。因此，两者各自接受的目标值序列非负、单调有界，因而收敛。实际计算中，AP检查共同投影梯度；VP-BNOMP检查相邻循环的归一化残差能量变化，并设置安全迭代上限。局部搜索所得位置还取决于初始化，这也是保留多个候选的原因。

\subsection{计算量及方法关系}\label{sec:complexity}
复杂度以保留频点数$F$、站数$L$、源数$K$和初始网格点数$G=|\mathcal G|$表示。记$I$为循环精修次数，$T$为连续目标或梯度的调用总数。单次联合系数求解需形成$K\times K$ Gram矩阵$\bm A_f^H\bm A_f$，并解相应线性方程，跨频主要计算量为$\mathcal O(F(LK^2+K^3))$。固定其他源后，投影可预计算，条件网格扫描约为$\mathcal O(FGLK)$。双源全组合搜索需检查$G^2$组位置，而逐源发现每次只扫描$G$点。

\begin{table}[pos=htbp]
\centering\small
\caption{主要计算量。$I_{\mathrm B}$为稀疏贝叶斯迭代数，$S$为分段数。}\label{tab:complexity}
\begin{tabularx}{\linewidth}{@{}lX@{}}\toprule
方法或运算 & 主要量级与可观测计数\\\midrule
公共DFT & $\mathcal O(LN\log N)$；原始采样点与保留频点数\\
MUSIC／MVDR & 统计约$\mathcal O(F_sSL^2+F_sL^3)$，谱扫描$\mathcal O(F_sGL^2)$；$F_s$为分段后频点数\\
MCJ-AP & 两次条件扫描加$M$个局部位置优化；约$\mathcal O(FGLK+T_MF(LK^2+K^3))$\\
VP-BNOMP & 逐源扫描与$I$轮循环精修；约$\mathcal O(KFGL+IKT_{\mathrm{loc}}FL)$，另计共同系数重估\\
多初值／GA & 由位置目标调用总数$T$决定，约$\mathcal O(TF(LK^2+K^3))$\\
宽带SBL & 稠密字典每轮约$\mathcal O(FGL^2+FL^3)$，累计随$I_{\mathrm B}$增长\\\bottomrule
\end{tabularx}
\end{table}
表中$T_M$为全部$M$个候选分支的连续目标与梯度调用总数，$T_{\mathrm{loc}}$为一次源位置精修的调用数，$F_s$为分段DFT保留的频点数。本文通过减少实际调用次数节省计算。局部单源系数按频点独立，式\eqref{eq:schur}的$H_{uu}$为分块对角矩阵，每块仅$2\times2$；逐频计算后累加得到$2\times2$位置曲率，无需构造完整$2F\times2F$矩阵。固定系数与消元更新的单次运算均随$F$线性增长，差别主要在所需步数与循环次数。实验因此同时报告扫描点数、目标/梯度调用数和完整耗时。

表\ref{tab:methodrelation}对齐本文与近邻方法的系数处理和位置更新。宽带最小二乘、变量投影与准Newton是共同工具，本文的改进分别落在条件候选的保留与比较、块NOMP的局部消元更新。两种流程服务于候选覆盖和更新效率这两个环节，实验分别检验其作用。

\begin{table}[pos=htbp]\centering\small
\caption{近邻方法与本文流程的关系。表中比较算法结构，不将不同观测模型下的原文性能直接排名。}\label{tab:methodrelation}
\begin{tabularx}{\linewidth}{@{}>{\raggedright\arraybackslash}p{0.17\linewidth}XX@{}}\toprule
方法 & 未知系数及位置求解 & 与本文的关系\\\midrule
宽带ML／AP\cite{ChenHudsonYao2002} & 消去未知波形；固定其余位置后作条件更新 & 提供联合目标与条件投影依据\\
二维多频BNOMP\cite{Yang2022} & 逐频源系数未知；固定当前系数求局部Newton方向，更新后重估，循环后共同LS & VP-BNOMP保留发现与循环时序，改用局部消元曲率和接受目标\\
QNOMP\cite{Zeng2025} & 已知导频信道观测；消去路径增益，联合BFGS更新时延和角度 & 提供消元后联合优化先例；本文估计逐源逐频波形，并直接优化位置\\
LAOLS\cite{Chatterjee2012} & 保留多个离散原子候选，向前补全支持集，比较最终拟合残差 & 提供延后候选决策的先例；MCJ-AP比较允许两源连续移动后的条件分支\\
MCJ-AP & 条件谱保留多个起点；每分支共同重估波形并联合精修位置 & 在最终拟合残差上比较候选，减少粗排序过早决策\\
VP-BNOMP & 每次局部位置试探都重估当前源系数；逐源循环后共同LS & 去除固定系数曲率中的耦合项，减少重复小步\\\bottomrule
\end{tabularx}\end{table}

\FloatBarrier
% END sections/methods


% BEGIN sections/experiments
\section{仿真实验}\label{sec:experiments}
\subsection{场景、比较方法与评价设置}\label{sec:baselines}
四站坐标分别为$(500,0)$、$(0,500)$、$(-500,0)$和$(0,-500)$\,m。以最大站间距离定义基线$R=\max_{l,v}\|\bm r_l-\bm r_v\|_2=1000$\,m。搜索区域为$[-1000,1000]^2$\,m$^2$，真实源位置为连续坐标，不要求落在初始网格上。两源采用独立BPSK符号，经根升余弦（RRC）脉冲成形；符号率$R_s=10$\,MHz、滚降系数$\alpha_{\mathrm{RRC}}=0.25$，对应名义带宽$B=(1+\alpha_{\mathrm{RRC}})R_s=12.5$\,MHz。表\ref{tab:scenario}汇总公共设置。

\begin{table}[pos=H]\centering\small
\caption{公共仿真参数。}\label{tab:scenario}
\begin{tabularx}{.96\linewidth}{@{}lX@{}}\toprule
参数 & 设置\\\midrule
接收站／源数 & 四个同步单天线站；已知$K=2$\\
传播 & 自由空间相对幅度$1/d$与基带时延；$c_0=299792458$\,m/s\\
波形 & 独立BPSK；RRC成形，符号率10\,MHz，滚降0.25\\
中心与名义带宽 & 两源中心均为0；带宽均为12.5\,MHz\\
采样与记录 & 50\,MHz；8192点；163.84\,$\mu$s\\
处理频带 & $[-7,7]$\,MHz；完整记录DFT保留2293个频点\\
噪声与信噪比 & 时域复噪声方差1；源B四站平均带内SNR为6\,dB\\
平均接收功率差 & 0、3、6、9、12、15\,dB；固定源B与噪声，仅缩放源A\\
初始网格 & 各方法统一20\,m，$101\times101=10201$点\\
候选间距 & 30\,m；AP联合版本分别保留1、2、5个候选\\
正式规模 & 20布局$\times$20随机实现$\times$6档功率差，共2400条\\
共同子集 & 各布局预先固定2次实现，保留全部功率差，共240条\\\bottomrule
\end{tabularx}
\end{table}

两源的符号、波形初相和定时独立。同一布局与随机实现的六档功率差共用波形和噪声，因而可以直接观察另一源增强带来的变化。源A为功率扫描源，源B为固定源；在正功率差下分别称强源和弱源，0\,dB时采用A／B标记。各方法只处理预先指定的公共频带$[-7,7]$\,MHz，在名义带宽两侧各留0.75\,MHz保护带，不提供逐源波形或频谱。筛选频点不重新缩放DFT，也不作逐站混合功率归一化。

正式比较采用20个布局，覆盖双源均在站内、A外B内、A内B外和双源均在站外四类构型，每类5个。站内指四站凸包内部，站外仍位于搜索区域内。几何生成排除距任一站小于150\,m、距中心小于100\,m、距站点对连线小于50\,m的位置，两源间距至少300\,m。位置按预定几何约束生成，不按定位结果筛选。图\ref{fig:formalgeometry}展示全部布局。

实验先在共同240条观测上比较八类方法，再将核心方法的随机重复从每布局2次增加到20次，形成2400条扩展验证。两组均覆盖全部20布局和六档功率差，区别在于随机重复次数；240条是2400条的预先固定子集，不额外增加样本总数。全面比较用于衡量各方法的可靠性与成本，扩展验证则用于检查核心方法的表现及收益来源。
\begin{figure}[pos=htbp]\centering
\includegraphics[width=.88\linewidth]{figures/formal_geometry.pdf}
\caption{20个源站布局，每类5个。黑色三角形为接收站，灰色菱形边界为站点凸包；同一面板内，相同形状的两个标记为一对源，红色为源A、蓝色为源B。0\,dB时两源平均接收功率相同。}\label{fig:formalgeometry}
\end{figure}
\FloatBarrier
图\ref{fig:formalgeometry}覆盖站内与站外的四类组合，使方法比较同时包含不同的源站位置关系。同一面板用五种形状区分五对源，颜色只区分源A与源B；各功率条件保持对应位置不变。

所有方法共享IQ观测、已知传播信息、搜索区域和20\,m初始网格，保留各自的连续精修。LS谱、MUSIC、MVDR、AP、块NOMP、SBL、多初值联合ML和GA联合ML的本场景准则列于附录\ref{app:baselines}；具体参数见表\ref{tab:parameters}。多初值ML采用去重的谱峰组合与Sobol空间覆盖起点，GA采用128个体，避免因重复初始化削弱搜索覆盖。MUSIC以同一记录的四个片段估计协方差，按扣除噪声后的观测功率加权；MVDR-known使用八个片段，相对对角加载为0.01。两者均只分析公共处理带内的频点，不将MVDR-known称为原版DPD-HR。
\par
% BEGIN tables/generated_parameters
\begin{table}[pos=H]\centering\footnotesize
\caption{正式比较采用的方法参数。连续位置优化不受20\,m网格量化。}\label{tab:parameters}
\begin{tabularx}{\linewidth}{@{}p{.18\linewidth}p{.27\linewidth}X@{}}\toprule
版本 & 核心操作 & 参数\\\midrule
LS谱 & 单源LS谱取两个峰 & 共同20 m网格；同谱连续精修\\
MUSIC & 噪声子空间宽带汇聚 & 4段；超噪声功率加权；已知K\\
MVDR-known & 加载协方差逆谱 & 8段；相对加载0.01\\
充分AP & 反复全域条件更新 & 候选5；循环上限6；正常连续精修\\
交替／联合AP & 相同分支，改变精修方式 & 候选1/2/5；每次局部上限100；共同投影梯度阈值$10^{-8}$\\
VP-BNOMP & 残差发现、消元Newton、共同LS & 循环上限10000；能量变化阈值$10^{-8}$；每次Newton最多5步；步长上限50 m；局部小步阈值0.01 m\\
固定系数BNOMP & 原文对应的固定系数更新 & 循环上限10000；能量阈值$10^{-8}$；局部小步阈值0.01 m\\
多初值联合ML & 同一联合LS目标 & 8个唯一初值：4谱峰+4 Sobol；L-BFGS-B\\
GA联合ML & 同一联合LS目标 & 种群128；最多200代；连续30代停滞停止\\
宽带SBL & 共享功率、离网字典、单点$\gamma$ & 最多300轮；每5轮离网更新；8个活动点；原生停止\\
\bottomrule\end{tabularx}\end{table}

% END tables/generated_parameters
\par
表\ref{tab:parameters}区分每种方法的搜索、精修和停止设置。各方法共享初始网格，但不强求相同迭代次数；AP按共同梯度精度比较，固定系数与消元块NOMP则保留相同原生循环停止规则。

% BEGIN tables/generated_environment
运行环境为Windows、Python 3.10.20、NumPy 2.2.6和PyTorch 2.8.0+cu129，GPU为NVIDIA GeForce RTX 4070 Ti SUPER，CPU计算线程设为20。响应与目标使用双精度实数／复数运算。

% END tables/generated_environment
\par

\FloatBarrier
\paragraph{定位评价}
设真实与估计集合分别为$X=\{\bm p_k\}$、$Y=\{\widehat{\bm p}_j\}$。每条观测先采用以平方欧氏距离为代价的匈牙利算法确定一对一匹配\cite{Kuhn1955}，再由同一组匹配误差计算RMSE、CDF、分位数和成功率：
\begin{equation}
 \pi^*=\arg\min_{\pi}\sum_{k=1}^{K}\|\bm p_k-\widehat{\bm p}_{\pi(k)}\|_2^2,
 \qquad e_{ik}=\|\bm p_{ik}-\widehat{\bm p}_{i,\pi^*(k)}\|_2 .\label{eq:assignment}
\end{equation}
其中$i$为观测编号，$k$为源编号，$\pi$为估计编号的排列，$e_{ik}$单位为m。对$n$条观测，RMSE和经验累计分布函数（CDF）为
\begin{equation}
 \operatorname{RMSE}_k=\sqrt{\frac1n\sum_{i=1}^{n}e_{ik}^2},\qquad
 \widehat F_k(t)=\frac1n\sum_{i=1}^{n}\mathbf 1(e_{ik}\le t).\label{eq:rmseCDF}
\end{equation}
式中$\mathbf 1$为示性函数，P95为第95百分位数，分位数采用相邻次序统计量的线性插值。RMSE和CDF使用全部样本，包括未达到成功门槛的样本。容许误差预先设为最大站点基线的2\%，即$r=0.02R=20$\,m，与初始网格步长分别定义。第$k$源成功率为$\widehat F_k(r)$，双源成功率为
\begin{equation}
 \widehat P_{\mathrm{both}}=\frac1n\sum_{i=1}^{n}
 \mathbf 1\!\left(\max\{e_{i1},e_{i2}\}\le r\right).
\end{equation}
同一估计不能对应两个真源。缺失或无效输出计为未成功，不从成功率分母中删除；本批正式结果均返回两个有限位置。误差CDF展示定位精度的整体分布，RMSE反映大误差的影响，双源成功率则衡量同时满足给定精度要求的比例。

\paragraph{成本与统计}
完整算法时间包含IQ输入后的DFT、频带筛选、传输、搜索、系数重估、精修及输出检查，GPU计时同步，公共预热单列。正式适配器内部的轨迹记录与最终目标、梯度检查均计入完整耗时；输入文件读取、结果落盘和离线统计另行记录。各方法独立承担搜索与初始化成本。目标、梯度、响应及循环调用分别统计，包含梯度计算的目标请求不与梯度次数相加。布局内同一实现的六档功率差保持配对；不确定性按构型类别、布局及随机实现分层重采样。核心比较和共同子集比较分别统计。

\subsection{八类方法的共同观测比较}\label{sec:commonformal}
首先比较各类方法在相同混合观测上的恢复能力与完整计算成本。表\ref{tab:formalcommon}使用各布局两次固定实现及全部六档功率差，共240条，所有方法逐条配对。VP-BNOMP采用同批运行的消元实现，各方法成本均来自这组共同观测。

% BEGIN tables/generated_formal_common
\begin{table}[pos=H]\centering\small
\caption{同一240条观测上的基线与提出方法比较。成功率括号内为成功数/总数；RMSE单位为m，时间均值、中位数及P95单位为s。}\label{tab:formalcommon}
\begin{tabular}{@{}lrrrrrr@{}}\toprule
方法 & 双源成功率/\% & 强源RMSE & 弱源RMSE & 时间均值 & 中位数 & P95\\\midrule
LS谱取峰 & 15.83 (38/240) & 191.982 & 845.184 & 0.2368 & 0.2181 & 0.3297\\
宽带MUSIC & 62.92 (151/240) & 93.515 & 646.892 & 0.2015 & 0.1893 & 0.2719\\
MVDR-known & 31.25 (75/240) & 320.684 & 804.082 & 0.1925 & 0.1807 & 0.2625\\
充分AP & 100.00 (240/240) & 0.876 & 1.596 & 8.9076 & 9.1483 & 12.2348\\
固定系数BNOMP & 95.00 (228/240) & 0.906 & 89.706 & 2.9416 & 1.4625 & 8.1989\\
VP-BNOMP & 95.00 (228/240) & 0.875 & 86.177 & 0.5907 & 0.3921 & 1.3971\\
多初值联合ML & 18.75 (45/240) & 256.061 & 776.935 & 1.2227 & 1.0868 & 2.5595\\
GA联合ML & 77.08 (185/240) & 279.056 & 501.231 & 8.0702 & 8.1102 & 9.0229\\
宽带SBL & 100.00 (240/240) & 1.096 & 3.743 & 44.1296 & 36.8237 & 97.7075\\
五候选联合AP & 100.00 (240/240) & 0.876 & 1.596 & 0.7081 & 0.6347 & 1.0232\\
\bottomrule\end{tabular}\end{table}

% END tables/generated_formal_common

在相同观测上，MCJ-AP以更低成本达到了充分AP和SBL的恢复水平。三者均全部成功，平均完整耗时分别为0.708\,s、8.908\,s和44.130\,s。固定系数BNOMP与VP-BNOMP均成功228条；消元版本的平均耗时由2.942\,s降至0.591\,s。两者保留相同发现和循环后共同LS结构，局部系数处理的差别在第\ref{sec:mechanismresults}节分析。
\begin{figure}[pos=htbp]\centering
\includegraphics[width=\linewidth]{figures/assignment_cost_20261019/common_error_cdf.pdf}
\caption{共同240条观测的强、弱源误差CDF。上行展示完整误差范围，横轴在1\,m内线性、其外对数；下行为小误差区间的线性放大。所有面板均使用全部240条作分母，竖线为20\,m参考线。}\label{fig:commonerrorcdf}
\end{figure}
\FloatBarrier
图\ref{fig:commonerrorcdf}将定位精度与恢复概率放在同一分布中观察。MCJ-AP和充分AP的曲线几乎重合，两源误差均集中在较小范围；SBL也恢复全部双源，但弱源误差分布更宽。VP-BNOMP在大部分观测上接近AP的精度，仍有一段远距离错误尾部。谱方法与多初值ML的低成功率同时伴随较大误差，并非仅由20\,m门槛附近的轻微偏差造成。

\begin{figure}[pos=htbp]\centering
\includegraphics[width=\linewidth]{figures/assignment_cost_20261019/common_time_cdf.pdf}
\caption{与表\ref{tab:formalcommon}同批的完整耗时CDF。左图覆盖全部耗时，横轴0.1\,s内线性、其外对数；右图以线性坐标放大早期完成区间。}\label{fig:commontimecdf}
\end{figure}
\FloatBarrier
图\ref{fig:commontimecdf}显示，谱方法计算较快，但应结合误差CDF判断其定位质量。MCJ-AP在保持高恢复率的同时，比充分AP和SBL更早完成；VP-BNOMP多数观测较快，少量长循环使其耗时尾部超过MCJ-AP。因此，中位耗时与P95分别补充了典型成本和慢样本成本，不能只看平均值。

\begin{figure}[pos=htbp]\centering
\includegraphics[width=\linewidth]{figures/formal_common.pdf}
\caption{共同240条观测的功率差扫描，每档40条。左侧比较谱方法，右侧比较逐源与联合求解方法；两个面板使用同一五候选AP参照。}\label{fig:formalcommon}
\end{figure}
\FloatBarrier

\begin{samepage}
图\ref{fig:formalcommon}显示，当前分段MUSIC与MVDR-known在功率差增大后明显退化，整体成功率分别为62.92\%和31.25\%；LS谱取峰为15.83\%。解释这组谱方法结果时，还需考虑分段统计和有限窗近似：附录\ref{app:segmentcheck}量化了后者的观测误差，其对定位结果的具体贡献尚未分离。多初值ML成功率为18.75\%，充分AP则达到100\%；二者使用同一联合目标，结果差别体现了初始化与搜索路径的重要性。GA成功率为77.08\%，平均耗时为8.070\,s。
\end{samepage}

MCJ-AP兼顾高恢复率与较低计算成本，VP-BNOMP则以较少计算取得95\%的双源成功率。SBL有226条按原生稳定条件停止、14条达到300轮上限。共同子集的停止状态见表\ref{tab:auxiliary}。

\FloatBarrier
\subsection{核心方法的2400条扩展验证}\label{sec:coreformal}
增加随机重复后，MCJ-AP仍保持较高的恢复率，VP-BNOMP保持较低的逐源更新成本。表\ref{tab:formalcore}汇总全部20布局、2400条观测；VP-BNOMP采用第\ref{sec:vpnomp}节的消元Newton实现。

% BEGIN tables/generated_formal_core
\begin{table}[pos=H]\centering\small
\caption{全部2400条观测的核心比较。成功率括号内为成功数/总数；RMSE单位为m，时间均值、中位数及P95单位为s。}\label{tab:formalcore}
\begin{tabular}{@{}lrrrrrr@{}}\toprule
方法 & 双源成功率/\% & 强源RMSE & 弱源RMSE & 时间均值 & 中位数 & P95\\\midrule
单候选联合AP & 95.00 (2280/2400) & 0.897 & 107.265 & 0.3958 & 0.3700 & 0.5104\\
双候选联合AP & 95.00 (2280/2400) & 0.897 & 103.698 & 0.4704 & 0.4329 & 0.6340\\
五候选联合AP & 99.00 (2376/2400) & 0.896 & 46.984 & 0.7180 & 0.6495 & 1.0267\\
VP-BNOMP & 94.92 (2278/2400) & 0.892 & 74.193 & 0.5855 & 0.3693 & 1.2470\\
\bottomrule\end{tabular}\end{table}

% END tables/generated_formal_core

五候选MCJ-AP成功2376条，将单候选的95\%成功率提高至99\%，平均耗时为0.718\,s。两者的强源RMSE均约0.897\,m，弱源恢复没有降低强源精度。五候选成功率的95\%分层重采样区间为$[96.83\%,100\%]$；相对双候选的成功率增加4个百分点，对应配对区间为$[0,12.17]$个百分点。恢复收益集中在本批困难布局，逐布局分布见第\ref{sec:timecases}节。VP-BNOMP成功2278条，成功率为94.92\%，平均耗时0.586\,s。该版本由单独完成的同输入验证计时，严格的消元更新效率比较采用第\ref{sec:mechanismresults}节的同批共同观测。

功率差既可能影响正确区域内的定位精度，也可能改变候选选择。图\ref{fig:formalcore}将本批诊断发现的布局4与其余19布局分开，分别观察这两种影响。两组均保留全部随机实现、功率条件和失败结果，总体成绩见表\ref{tab:formalcore}。
\begin{figure}[pos=htbp]\centering
\includegraphics[width=.94\linewidth]{figures/formal_core.pdf}
\caption{扩展验证中功率差的分组结果。左列为其余19布局，共2280条、每档380条；右列为困难布局4，共120条、每档20条。上行为双源成功率，下行为弱源RMSE。分组依据本批构型诊断，全部2400条仍计入总体统计。}\label{fig:formalcore}
\end{figure}
\FloatBarrier
其余19布局中，两种AP均全部成功，MCJ-AP的弱源RMSE随功率差从1.31\,m增至1.93\,m，显示强源增强后弱源定位精度逐步下降。VP-BNOMP在这一组有两条15\,dB失败，因此该档RMSE仍受到长尾影响。右列单独呈现候选恢复困难，具体失败分布和候选机制在第\ref{sec:timecases}节集中分析。

\begin{figure}[pos=htbp]\centering
\includegraphics[width=\linewidth]{figures/assignment_cost_20261019/core_error_cdf.pdf}
\caption{全部2400条核心观测的强、弱源误差CDF。上行横轴在1\,m内线性、其外对数，下行线性放大小误差区间；每条曲线保留全部样本，竖线标出20\,m。}\label{fig:formalcdfcost}
\end{figure}
\FloatBarrier
图\ref{fig:formalcdfcost}中，四种版本的强源CDF几乎重合。弱源小误差区间内，VP-BNOMP与单、双候选AP接近；远端尾部不同，使VP-BNOMP的总体RMSE低于这两种AP。五候选MCJ-AP则恢复更多弱源，CDF在20\,m附近达到更高水平，其弱源P95为3.85\,m、RMSE为46.98\,m。CDF与RMSE共同表明：多数案例的精度、未恢复比例和严重错误大小是不同的性能方面。

\FloatBarrier
\subsection{候选保留提高可靠性，联合精修减少计算}\label{sec:apresults}
图\ref{fig:formalablation}分别检验候选数量与位置更新方式：前者决定搜索覆盖，后者决定处理同一组候选所需的计算。
\begin{figure}[pos=htbp]\centering
\includegraphics[width=\linewidth]{figures/formal_ablation.pdf}
\caption{AP的候选与精修消融，四面板单行排列。左两图为原2400条主批的成功率与耗时，右两图为同一2400条输入上的五候选配对耗时与梯度计算。ALT为交替精修，joint为联合精修；左、右两批分别计时，不跨批计算加速比。}\label{fig:formalablation}
\end{figure}
\FloatBarrier
左两图显示，候选数为1和2的联合AP均成功2280条，五候选增至2376条。第二候选尚未覆盖本批困难区域，保留更多有效起点才带来恢复。右两图显示，在候选数相同的条件下，联合精修同时减少实际耗时和梯度计算。

\paragraph{候选数量与完整计算成本}
为直接比较候选数量与逐源更新的实际取舍，在共同240条观测上交错运行候选数$M=1,2,3,4,5$的联合AP和VP-BNOMP，共1440次求解。各版本独立完成预处理、候选搜索、精修和输出检查，配置及原生停止条件保持不变。以下耗时均来自这一同批重测，不与表\ref{tab:formalcommon}的历史墙钟拼接计算加速比。三、四候选的位置与五分支前缀择优结果逐例一致。

\begin{figure}[pos=htbp]\centering
\includegraphics[width=\linewidth]{figures/assignment_cost_20261019/candidate_time_cdf.pdf}
\caption{同批240条观测的AP候选数与VP-BNOMP完整耗时CDF。每条曲线包含240条观测的完整求解耗时；左图0.1\,s内线性、其外对数，右图线性放大0--1.5\,s。}\label{fig:candidatetimecdf}
\end{figure}
\FloatBarrier
图\ref{fig:candidatetimecdf}中，增加AP候选数使耗时分布逐步右移。VP-BNOMP的中位耗时为0.3674\,s，接近单候选AP的0.3514\,s，但其耗时P95为1.3147\,s，高于五候选AP的0.7201\,s。两种流程的成本差别主要体现在分支数量和少量长循环，而不只是平均速度。

\begin{figure}[pos=htbp]\centering
\includegraphics[width=\linewidth]{figures/assignment_cost_20261019/candidate_tradeoff.pdf}
\caption{同一240条观测的可靠性、误差与平均完整耗时。误差棒为固定20布局内按随机实现配对重采样的95\%区间，同一实现的六档功率差保持在一起；区间不表示对任意几何总体的推断。}\label{fig:candidatetradeoff}
\end{figure}
\FloatBarrier
图\ref{fig:candidatetradeoff}中，$M=1$至$4$均成功228条，平均耗时从0.3556\,s增至0.5420\,s，弱源RMSE约111.2\,m；VP-BNOMP同样成功228条，以0.5474\,s平均耗时取得86.18\,m弱源RMSE。它与四候选AP的平均耗时接近，但部分错误位置的偏差更小。五候选AP成功240条，弱源RMSE降至1.60\,m，平均耗时0.6079\,s。因而，单候选AP强调低成本，VP-BNOMP降低了本批小候选AP的严重错误幅度，五候选AP提供更高恢复率；这一取舍不能概括为VP-BNOMP总体最快。

在双方均成功的228条观测中，VP-BNOMP有197条比五候选AP更快；五候选AP则额外恢复其余12条。三、四候选在全部2400条已保存分支上的离线位置消融分别成功2281和2289条，说明共同子集没有包含所有候选恢复变化；该全批消融只评价位置，不将240条耗时外推为2400条实测成本。

单候选配对先说明了联合精修的省时作用：相同成功数和接近的RMSE下，平均耗时从0.439\,s降至0.396\,s，减少9.84\%。平均耗时差（联合减交替）的95\%分层配对区间为$[-0.0496,-0.0367]$\,s；目标计算从39.76次降至28.69次，梯度计算从37.76次降至26.69次。

为确认省时不是降低求解精度造成的，表\ref{tab:fivepair}单独列出五候选配对的精度和尾部指标。两种版本对每条观测使用相同的五个初始候选，在同一完整目标和投影梯度阈值下停止；运行顺序交错，共2400对、4800次求解。

% BEGIN tables/generated_formal_pair
\begin{table}[pos=H]\centering\small
\caption{相同五候选配对的定位精度，2400对；计算成本见图\ref{fig:formalablation}。}\label{tab:fivepair}
\begin{tabular}{@{}lrrrr@{}}\toprule
精修方式 & 双源成功 & 强源RMSE/m & 弱源RMSE/m & 弱源P95/m\\\midrule
交替 & 2376 & 0.896 & 46.983 & 3.846\\
联合 & 2376 & 0.896 & 46.984 & 3.845\\
\bottomrule\end{tabular}\end{table}

% END tables/generated_formal_pair

表\ref{tab:fivepair}中，两者均成功2376条，成功／失败身份逐例一致；强源RMSE约0.896\,m，弱源P95约3.85\,m，弱源RMSE也几乎相同。联合更新保持了交替更新的定位质量。图\ref{fig:formalablation}右两图显示，平均完整耗时从1.0840\,s降至0.7441\,s，减少31.36\%；耗时差的95\%分层配对区间为$[-0.3691,-0.3121]$\,s。20个布局的平均耗时均减少，逐布局节省11.42\%--48.04\%，省时作用不依赖单一构型。

省时与计算量减少一致：平均联合目标计算从204.21次降至120.36次，梯度计算从198.21次降至114.36次，后者减少约42.30\%。两种版本选中分支均满足共同梯度条件，全部候选中有两次交替精修触及安全迭代上限，其结果仍计入统计。由此，多候选保留了困难弱源的有效起点，联合精修则以更少局部计算处理这些候选，两者共同构成MCJ-AP的可靠性与效率收益。

\FloatBarrier
\subsection{波形消元更新的机制与效率}\label{sec:mechanismresults}
式\eqref{eq:schur}给出了固定系数与消元曲率之间的关系。位置改变时，波形系数可以补偿一部分响应变化，固定系数曲率没有扣除这部分耦合，可能产生反复的小步。局部消元同时改变更新方向和试探位置的目标计算，因此应与保留相同外层时序的参照配对，区分局部更新收益与整套实现差别。

为区分更新机制、停止尺度和数值保护的影响，采用四类开发构型、每类一个随机实现及0/15\,dB两档功率差，共8条观测作受控诊断。表\ref{tab:mechanism}中各变体只改变所列因素，候选发现和循环后共同LS保持一致。这组小规模对照用于解释机制，正式性能由后续共同240条及完整2400条结果评价。

\begin{table}[pos=H]\centering\small
\caption{局部更新的受控机制诊断。各版本使用相同8条开发观测；耗时包含必要预处理与求解，导数调用为每例均值。}\label{tab:mechanism}
\begin{tabular}{@{}lrrr@{}}\toprule
局部更新设置 & 双源成功 & 耗时/s & 导数调用\\\midrule
固定系数，能量容差$10^{-8}$ & 8/8 & 5.478 & 2679.5\\
仅放宽能量容差至$10^{-7}$ & 7/8 & 2.318 & 1133.9\\
仅放宽能量容差至$10^{-6}$ & 6/8 & 0.971 & 467.4\\
仅取消局部小步长退出 & 8/8 & 8.011 & 4758.4\\
局部波形消元，能量容差$10^{-8}$ & 8/8 & 0.436 & 70.4\\\bottomrule
\end{tabular}\end{table}

固定系数参照的21436次局部更新全部被接受，其中83.67\%的步长小于0.01\,m；Hessian正定修正、步长封顶、回溯缩步与边界裁剪均未触发。小步因此来自更新方向本身。放宽停止虽减少计算，却使部分弱源误差超过20\,m；取消小步长退出则增加了每轮内部计算。局部消元保留原能量容差，将平均导数调用由2679.5次降至70.4次。

相同位置与最小二乘系数下的曲率核对进一步解释了这一差别。在两个站外诊断位置，未施加数值保护的固定系数Newton试探步分别约0.0138\,m和0.0111\,m，消元试探步分别约11.25\,m和32.26\,m；这些试探方向用于比较曲率，实际接受步仍由回溯确定。四个保存位置的解析消元Hessian与自动微分结果的最大差约$1.4\times10^{-18}$。这与式\eqref{eq:schur}一致：允许波形随位置调整后，消元曲率去除了可由系数补偿的响应变化，减少了固定系数更新中的重复小步。

在同批共同240条上，固定系数与消元版本都成功228条，平均完整耗时分别为2.9416\,s和0.5907\,s，消元版本减少约79.92\%的计算时间，见表\ref{tab:formalcommon}。这一对照保持残差发现、逐源循环和循环后共同LS的时序，局部更新将固定系数Newton及其回溯目标改为式\eqref{eq:localprofile}的消元形式。结合式\eqref{eq:schur}，它直接检验了位置--波形耦合处理对计算效率的作用。

在全部2400条上，VP-BNOMP成功2278条，整体精度与成本见表\ref{tab:formalcore}。成功观测中的弱源RMSE约1.52\,m，表中的总体RMSE则同时计入未恢复观测。MCJ-AP通过分支保留扩大有效候选覆盖，VP-BNOMP通过局部消元降低逐源更新成本，两者的贡献分别由候选消融与固定系数对照检验。

VP-BNOMP全部观测均按相邻循环的归一化残差能量变化不超过$10^{-8}$停止，平均记录约51.03轮，没有达到10000轮安全上限。该上限仅防止异常长循环，不代表实际需要接近此数量的迭代。每次局部Newton最多5步，位置试探上限50\,m，接受小步小于0.01\,m时结束本次局部调用。推导、伪代码、参数表及上述统计均对应同一消元实现。

\subsection{求解时间轨迹与困难构型的候选恢复}\label{sec:timecases}
图\ref{fig:formaltime}重放共同240条观测上已经保存的输出。尚无双源输出的观测按未成功计；RMSE仅在全部观测有输出后显示。
\begin{figure}[pos=htbp]\centering
\includegraphics[width=\linewidth]{figures/formal_time.pdf}
\caption{共同240条观测的保存输出轨迹。上行使用0.25--300\,s对数时间轴，下行以线性时间轴放大0--3\,s；三列依次为双源成功率、完整输出比例和两源合并RMSE，后者为$\sqrt{(nK)^{-1}\sum_{i,k}e_{ik}^2}$。正常停止后保持最终值，RMSE纵轴仍用对数刻度。连接线表示保存点之间的变化，不额外规定算法停止时刻。}\label{fig:formaltime}
\end{figure}
\FloatBarrier

五候选联合AP、充分AP和SBL首次达到全部成功的保存点分别为0.75\,s、1.5\,s和30\,s。取得高质量结果的时间与完成全部计算的时间需要分别观察。上行对数轴便于比较亚秒至数百秒的差异，下行则分辨早期输出。VP-BNOMP在本次登记轨迹中只发布最终完整双源结果，其曲线不代表内部每次循环均已有对外输出。

\paragraph{困难构型下的候选覆盖}
图\ref{fig:formallayouts}给出全部20布局的失败分布，用于定位总体长尾的来源。
\begin{figure}[pos=htbp]\centering
\includegraphics[width=\linewidth]{figures/formal_layouts.pdf}
\caption{各布局的双源失败数，每布局120条。四个面板对应四类构型，零表示该布局全部成功；不同标记在纵向稍作错开，便于辨认重合结果。}\label{fig:formallayouts}
\end{figure}
\FloatBarrier
单候选AP在布局4的120条观测上全部失败，五候选恢复其中96条，其余19布局没有成功转失败。VP-BNOMP在布局4全部失败，另外两条失败位于其余布局的15\,dB条件。多候选的可靠性收益因而集中在本批困难构型，下面通过候选位置和精修目标解释这一恢复过程。

MCJ-AP每档400条观测中有3--6条失败，均来自布局4，并贡献各档超过99.7\%的弱源平方误差。该布局内候选选择的变化主导了总体RMSE折线的非单调起伏；其余19布局的曲线则更清楚地显示功率差对定位精度的影响。同一组布局在六档功率差中保持不变，这种起伏并非各档采用了不同构型。

图\ref{fig:formalcandidates}给出同一困难布局的恢复与未恢复案例，以对照候选覆盖和最终选择。
\begin{figure}[pos=htbp]\centering
\includegraphics[width=\linewidth]{figures/formal_candidates.pdf}
\caption{同一困难构型的两个0\,dB事后机制案例。左列为单源LS谱，中列为条件谱，右列为候选精修前后残差；上行为恢复案例，下行为未恢复案例。红色和白色星号分别标示源A、源B，空心圆为条件候选。谱图采用对数色标，右列连接线仅连接候选编号，不表示实际迭代路径。}\label{fig:formalcandidates}
\end{figure}
\FloatBarrier
上行中，前四个候选都进入远端错误区域，第五候选则接近真实源B。联合精修后，第五分支得到更低残差，源B误差约0.09\,m。下行的五个候选均落在远端高分带，最近起点仍距源B约397\,m，局部精修未能恢复。这表明粗谱最高分与精修后的最佳位置并不等价；多候选增加有效起点，最终拟合负责择优，两者共同提高恢复能力。若候选全部落在错误区域，局部精修本身不能补足全域覆盖。

第\ref{sec:complexity}节的复杂度分析与实际调用量相对应：单／五候选AP均只作两次全域网格扫描，增加分支主要增加连续目标和梯度计算；固定系数与消元NOMP则需比较局部导数调用和循环成本。这些计数用于解释候选覆盖与局部更新各自带来的开销和收益。

% END sections/experiments

\FloatBarrier

% BEGIN sections/conclusion
\section{结论}\label{sec:conclusion}
本文针对未知宽带强弱重叠源，将波形消元与位置搜索结合，形成多候选联合AP和波形消元块NOMP。前者保留可能的起点，再联合调整两源位置；后者在位置试探时重估波形，减少固定系数带来的重复小步。

2400条观测中，五候选MCJ-AP达到99\%的双源20\,m成功率，相对单候选恢复96条，平均耗时0.718\,s。相同五候选配对中，联合精修保持定位表现，将平均耗时从1.084\,s降至0.744\,s，减少31.36\%，梯度计算减少42.30\%。构型分组与候选案例表明，多候选改善困难区域的恢复，联合精修降低分支处理成本。

VP-BNOMP在全部2400条上达到94.92\%成功率，平均耗时0.586\,s。同批共同240条中，消元更新保持固定系数版本的成功数，将平均耗时从2.942\,s降至0.591\,s，与曲率分析相符。两种流程分别强调候选覆盖与局部更新效率；后续可扩展至未知源数和更一般的传播条件。

% END sections/conclusion

\FloatBarrier
\appendix

% BEGIN sections/appendix
\section{比较方法在本模型下的准则}\label{app:baselines}
\subsection{单源谱、MUSIC及MVDR-known}
LS取峰使用式\eqref{eq:single}，每个位置分别按单源模型评分。MUSIC和MVDR则需要从观测估计空间统计量。为此，将同一记录划分为$S$个不重叠片段，各片段长度为$N_s=N/S$。令$\bm z_{bf}\in\C^L$表示第$b$段在频率$f$处的多站DFT观测，$b=1,\ldots,S$。各段采用共同的频率轴，并在该轴上选择相同物理频带；样本协方差估计为
\begin{equation}
 \widehat{\bm C}_f=S^{-1}\sum_{b=1}^S\bm z_{bf}\bm z_{bf}^H .
\end{equation}
同一记录内的位置、传播关系和信号统计保持不变，各片段DFT用于构成同频快拍。零均值模型下，相关矩阵等于协方差矩阵，故上式不另减样本均值。完整记录与片段DFT分别按$1/N$和$1/N_s$归一化，对应噪声方差为$\sigma_t^2/N$和$\sigma_t^2/N_s$。分段提供更多统计快拍，同时降低频率分辨率、缩短单段时长。

MUSIC利用信号响应与噪声子空间的正交关系定位\cite{WeissAmar2005}。其通常要求$K<L$、源频域系数的协方差满秩、不同源的跨站响应线性独立，并采用空间白噪声模型。在两源均有非零功率的频点，独立源的总体协方差满足相应源秩条件；有限记录下的子空间估计精度还与功率差、有效快拍数及频点信号强度有关。本文采用逐频响应与宽带评分汇聚。

设$\bm E_{n,f}\in\C^{L\times(L-K)}$的列为$\widehat{\bm C}_f$最小$L-K$个特征值对应的单位正交特征向量。用已知源数$K$划分子空间，本文MUSIC适配版使用
\begin{align}
 S_{\mathrm{MU}}(\bm p)&=\left[\sum_f w_f
 \bm q_f(\bm p)^H\bm E_{n,f}\bm E_{n,f}^H\bm q_f(\bm p)\right]^{-1},\label{eq:music}\\
 \widetilde w_f&=\max\{\tr(\widehat{\bm C}_f)-L\sigma_t^2/N_s,0\},\quad
 w_f=\widetilde w_f/\sum_v\widetilde w_v .\nonumber
\end{align}
式中$\tr$为矩阵迹，$\widetilde w_f$为扣除期望噪声后的非负观测功率，$w_f$为归一化权重，$v$遍历保留频点。所有权重为零时使用均匀权重。该宽带汇聚使信号能量较高的频点发挥更大作用；带外频点不参与分解。

MVDR在保持候选响应不失真的条件下，最小化输出功率\cite{Capon1969}。令$\bm h_f\in\C^L$为复数合成权重，对加载后的协方差解
\begin{equation}
 \min_{\bm h_f}\ \bm h_f^H\bm C_{f,\delta}\bm h_f,
 \qquad \bm h_f^H\bm q_f(\bm p)=1.\label{eq:mvdrConstraint}
\end{equation}
最优输出功率为$[\bm q_f^H\bm C_{f,\delta}^{-1}\bm q_f]^{-1}$。将这一功率在保留频点上平均，得到MVDR-known评分：
\begin{align}
 \bm C_{f,\delta}&=\widehat{\bm C}_f+
 \delta\,\tr(\widehat{\bm C}_f)\bm I/L,\nonumber\\
 S_{\mathrm{MV}}(\bm p)&=F_s^{-1}\sum_f
 [\bm q_f(\bm p)^H\bm C_{f,\delta}^{-1}\bm q_f(\bm p)]^{-1} .\label{eq:mvdr}
\end{align}
这里$\delta>0$为无量纲相对加载量，加载项是平均对角功率$\tr(\widehat{\bm C}_f)/L$的$\delta$倍，用于缓解有限快拍下求逆不稳定的问题；$F_s$为片段频率轴上保留的频点数，区别于采样率$f_s$。约束中的响应和评分均使用同一单位范数向量$\bm q_f$。“known”表示采用已知相对传播响应；本宽带汇聚式与DPD-HR的原公式分别命名\cite{TirerWeiss2016}。

\subsection{有限分段的模型一致性}\label{app:segmentcheck}
为检查短窗近似，使用四类开发构型的0/15\,dB观测，共8条，直接读取仿真保存的无噪声源分量。在真实位置处构造$\bm A_f$，计算干净混合观测的最小二乘残差；这一步只核对观测模型，不搜索位置。对$S$段、每段$N_s=N/S$点，定义
\begin{equation}
 E_{\mathrm{win}}=\sum_{b=1}^{S}\sum_{f\in\mathcal F_s}
 \|\bm\Pi_{\bm A_f}^{\perp}\bm z_{bf}^{(0)}\|^2,\qquad
 E_{\mathrm n}=SF_s(L-K)\sigma_t^2/N_s ,
 \label{eq:windowdiagnostic}
\end{equation}
其中$\bm z_{bf}^{(0)}$为无噪声分段DFT，$\mathcal F_s$为该段长下保留的频点集合，$F_s=|\mathcal F_s|$。$E_{\mathrm n}$是相同投影后的白噪声期望能量；各频点$\bm A_f$均经检查满列秩。表\ref{tab:windowdiagnostic}使用各自段长、频率轴和DFT归一化，传播常数与原仿真一致。

\begin{table}[pos=H]\centering\small
\caption{8条开发观测的有限窗模型核验。占比为$E_{\mathrm{win}}$除以保留频点的干净混合观测能量，表中列出观测间范围。}\label{tab:windowdiagnostic}
\begin{tabular}{@{}lrrr@{}}\toprule
分段数 & 单段时长/$\mu$s & 残差占比/\% & $E_{\mathrm{win}}/E_{\mathrm n}$\\\midrule
1 & 163.84 & 0.209--0.382 & 0.035--0.882\\
4 & 40.96 & 0.756--1.387 & 0.126--3.210\\
8 & 20.48 & 1.529--2.854 & 0.238--6.611\\\bottomrule
\end{tabular}\end{table}
分段越短，截取共同观测窗造成的跨站频域偏差越明显。15\,dB的四条观测中，4段和8段的模型残差均超过预计投影后噪声能量。因此，短窗模型误差是解释本组谱方法表现时需要考虑的因素；上述残差本身并不等同于定位误差。核心AP配对与BNOMP局部更新对照使用相同完整记录，不改变这一受控比较关系。

\subsection{充分AP、多初值联合ML与GA}
充分AP对各源反复进行全域条件扫描与连续精修，以式\eqref{eq:profile}判断更新\cite{ChenHudsonYao2002}。它与MCJ-AP优化同一目标，但进行更多初值处理和全域核查。实验同时报告保存的早期解与最终停止时间，区分取得结果和完成核查的成本。

多初值联合ML从四个唯一的谱峰组合和四个Sobol低差异序列给出的空间覆盖组合出发，以有界L-BFGS-B最小化$J$，输出最小残差解。候选组合的无序重复在初始化阶段删除。遗传算法（GA）将两个源的四个坐标编码为一个实数个体，在相同目标$J$上进行锦标赛选择、模拟二进制交叉和多项式变异\cite{DebAgrawal1995}，保留4个精英，最后作联合连续精修。锦标赛规模为3，交叉概率为0.9、分布指数为15，每维变异概率为0.25、分布指数为20。这些参数控制个体选择与扰动范围；总体成本由种群规模、代数和实际目标调用共同决定。

\subsection{多频共享功率SBL}
稀疏贝叶斯学习（SBL）在离散位置字典上估计稀疏源功率。令$\bm g_j\in\Omega$为第$j$个字典位置，$\bm D_f=[\bm q_f(\bm g_1),\ldots,\bm q_f(\bm g_G)]\in\C^{L\times G}$。逐频系数$\bm z_f\in\C^G$采用零均值复高斯先验，$\bm z_f\sim\mathcal{CN}(0,\bm\Gamma)$，其中$\bm\Gamma=\diag(\bm\gamma)$，$\bm\gamma\in\R_{\ge0}^G$为跨频共享的非负功率超参数\cite{Nannuru2019}。各频系数按该先验独立建模，位置支持通过共同的$\bm\gamma$关联。观测协方差、系数后验均值及待最小化的负对数边际似然分别为
\begin{align}
 \bm C_f&=\sigma_f^2\bm I+\bm D_f\bm\Gamma\bm D_f^H,\label{eq:sblCov}\\
 \bm\mu_f&=\bm\Gamma\bm D_f^H\bm C_f^{-1}\bm x_f,\nonumber\\
 \mathcal E(\bm\gamma)&=F^{-1}\sum_f[\log\det\bm C_f+\bm x_f^H\bm C_f^{-1}\bm x_f] .\label{eq:sblEvidence}
\end{align}
其中$\sigma_f^2=\sigma_t^2/N$，$\bm C_f\in\C^{L\times L}$，$\bm\mu_f\in\C^G$为后验均值；$\log\det$为自然对数行列式，$\mathcal E$省略了与超参数无关的常数。对每个频点采用同一功率先验时，当前实现的固定点更新为
\begin{equation}
 \gamma_j^+=\gamma_j\left[
 \frac{\sum_f|\bm d_{fj}^H\bm C_f^{-1}\bm x_f|^2}
 {\sum_f\bm d_{fj}^H\bm C_f^{-1}\bm d_{fj}}
 \right]^{1/2} .\label{eq:sblUpdate}
\end{equation}
这里$\bm d_{fj}$为$\bm D_f$的第$j$列，$\gamma_j^+$为更新后的超参数，指数$1/2$表示采用半步更新。每次更新通过回溯检查$\mathcal E$是否下降。每5轮对8个活动字典点作离网优化，即允许高权重点离开初始网格以减小量化误差。每维位移限制为20\,m，坐标尺度为1000\,m，L-BFGS-B最多20轮，最终从单点$\gamma_j$取峰。停止检查同样每5轮执行一次：负对数边际似然的归一化变化小于$10^{-5}$，$\gamma$的L1相对变化小于1\%，输出位置和活动点最大移动均小于0.05\,m；连续两次满足时停止，否则最多300轮。这些是优化稳定性条件，0.05\,m不是定位成功门槛。谱峰分组候选不并入本文提出方法。

\subsection{局部精修与停止参数}
AP的交替和联合版本在归一化目标$J$上使用相同的投影梯度阈值$10^{-8}$，目标差停止项设为零。投影梯度表示考虑搜索边界后的可行下降方向：在内部等于普通梯度，在边界排除指向区域外的分量。每次L-BFGS-B最多100轮，线搜索最多25次；交替版本最多12个完整位置扫描。初始化的单源精修采用目标相对容差$10^{-12}$。候选先在$3\times3$网格邻域中取局部极大值，再按30\,m间隔去重。连续精修不受初始20\,m网格步长限制，候选间隔也不等于成功门槛。

VP-BNOMP每次局部Newton最多5步，Hessian最小特征值下限为$10^{-9}$，位置步长上限为50\,m。按$1,1/2,1/4,1/8,1/16$依次回溯，每个试探位置重新估计信号系数，接受局部残差不增的更新；接受步长小于0.01\,m时结束本次局部调用。相邻完整逐源循环的归一化残差能量差不超过$10^{-8}$时停止，安全上限为10000轮。固定系数BNOMP使用相同数值保护与循环停止规则，但回溯时保持当前系数不变。

\subsection{固定系数BNOMP对应版与消元对照}
原文对应版根据给定系数$q_f^{(0)}$使用
\begin{equation}
 L_{\mathrm{FC}}(\bm p)=E_x^{-1}\sum_f
 \|\bm y_f-\bm a_f(\bm p)q_f^{(0)}\|^2 .\label{eq:fixedcoef}
\end{equation}
其中$\bm y_f=\bm y_{j,f}$为一次局部优化内固定的观测。每次局部调用先在当前位置计算条件最小二乘系数$q_f^{(0)}$，固定该系数求Newton方向；接受新位置后再次重估系数，供下一步求导使用。这使固定系数与消元曲率可在同一位置、同一最小二乘系数处比较。本文以二维多频块NOMP\cite{Yang2022}的固定系数位置方向、逐源循环和循环后共同LS为参照，将声学响应替换为四站时延--幅度响应。

原文的更新时序与本实现的数值保护分别说明：Hessian下限、步长限制、边界裁剪和回溯规则是本文补充的稳定化设置。参照版本在回溯中固定当前系数；VP-BNOMP则在每个试探位置重估系数。原始NOMP在固定增益曲率求方向后，以重估增益后的评分检查是否接受\cite{Mamandipoor2016}，因此固定系数曲率并不限定必须采用固定系数接受目标。本文的效率对照评价整套局部消元更新，曲率的作用另由第\ref{sec:mechanismresults}节的受控诊断说明。

本文VP-BNOMP采用受控消元版本：将式\eqref{eq:fixedcoef}替换为式\eqref{eq:localprofile}，并使用对应的消元曲率及回溯目标，保留原参照的发现和循环后共同LS时序。因此，固定系数与消元版本的比较对应同一外层结构中的局部位置--波形耦合处理。

\subsection{停止状态}
表\ref{tab:auxiliary}与正文共同子集表使用同一240条观测。达到迭代上限的结果仍按最终输出评价，其余停止原因及具体配置随可复现材料保存。

% BEGIN tables/generated_formal_auxiliary
\begin{table}[pos=H]\centering\small
\caption{共同240条观测中达到安全迭代上限的次数。}\label{tab:auxiliary}
\begin{tabular}{@{}lr@{}}\toprule
方法 & 达到迭代上限\\\midrule
LS谱取峰 & 0\\
宽带MUSIC & 0\\
MVDR-known & 0\\
充分AP & 0\\
固定系数BNOMP & 0\\
VP-BNOMP & 0\\
多初值联合ML & 0\\
GA联合ML & 208\\
宽带SBL & 14\\
五候选联合AP & 0\\
\bottomrule\end{tabular}\end{table}

% END tables/generated_formal_auxiliary

表\ref{tab:auxiliary}显示，充分AP、MCJ-AP和VP-BNOMP在共同子集中均按各自原生条件停止；GA与SBL分别有208条和14条达到安全上限。SBL的耗时P95约97.7\,s，反映了较慢样本的额外计算成本。

\section{消元局部导数与实现核验}\label{app:derivatives}
为计算Newton曲率，需要响应对位置的二阶导数。固定频率$f$，令$\kappa=2\pi f/c_0$，单位为rad/m；$d_l=\|\bm p-\bm r_l\|_2$、$\bm v_l=(\bm p-\bm r_l)/d_l$沿用正文定义。写$a_l=\exp(\ell_l)$，其中$\ell_l=-\log d_l-\mathrm j\kappa d_l$是这一响应的连续指数表示。在观测模型一节所述距离保护区外，$d_l>0$，有
\begin{align}
 \nabla\ell_l&=(-d_l^{-1}-\mathrm j\kappa)\bm v_l,\\
 \nabla^2\ell_l&=d_l^{-2}\bm v_l\bm v_l^T
 +\frac{-d_l^{-1}-\mathrm j\kappa}{d_l}
 (\bm I-\bm v_l\bm v_l^T),\\
 \nabla^2 a_l&=a_l[\nabla^2\ell_l+(\nabla\ell_l)(\nabla\ell_l)^T].\label{eq:responseHessian}
\end{align}
这里$\nabla$和$\nabla^2$分别对二维位置求梯度和Hessian，$\bm I$为$2\times2$单位矩阵。式\eqref{eq:responseHessian}同时包含幅度与时延的二阶变化；转置$T$没有取复共轭，因为此处求的是复响应对实坐标的普通导数。

对固定的$\bm y_f$，将系数$q_f$拆为实部和虚部，按$\bm u_f=[\Rea q_f,\operatorname{Im}q_f]^T\in\R^2$排列。该频点损失对$\bm u_f$的Hessian为
\begin{equation}
 H_{uu,f}=2\|\bm a_f\|^2\bm I_2/E_x .
\end{equation}
各频点的系数块相互独立，因此逐频计算混合导数，再累加式\eqref{eq:schur}的修正项，即得到$\phi_j$的$2\times2$实Hessian。实现中以直接最小二乘和小矩阵运算核对投影公式，以有限差分核对位置梯度与Hessian；定位性能另由仿真实验评价。

% END sections/appendix

\FloatBarrier
\bibliographystyle{\PaperRoot template/cas-model2-names}
\bibliography{\PaperRoot references/references}
\end{document}
